\documentclass[11pt,a4paper]{article}

\usepackage[T1]{fontenc}
\usepackage[utf8]{inputenc}
\usepackage{lmodern}
\usepackage[ngerman]{babel}
\usepackage[a4paper,margin=2.2cm]{geometry}
\usepackage{array}
\usepackage{longtable}
\usepackage{tabularx}
\usepackage{xcolor}
\usepackage{hyperref}

\setlength{\parindent}{0pt}
\setlength{\parskip}{0.5em}
\renewcommand{\arraystretch}{1.35}

\newcommand{\answerline}[1]{\rule{#1}{0.4pt}}
\newcommand{\term}[1]{\fcolorbox{black!20}{teal!8}{\strut #1}}

\title{\textbf{Aufgabenblatt: Protokolle und Schnittstellen}}
\author{Modul 321 -- Lektion 08}
\date{}

\begin{document}

\maketitle

\textbf{Name:} \answerline{7cm} \hfill \textbf{Datum:} \answerline{3.5cm}

\vspace{0.8em}

\fcolorbox{teal!40}{teal!6}{
  \parbox{\dimexpr\linewidth-2\fboxsep-2\fboxrule-6pt\relax}{
    \textbf{Ziel der Lektion:} Sie können wichtige Protokolle und Schnittstellenstile
    benennen, deren Aufgaben erklären und für einfache Szenarien sinnvoll auswählen.
  }
}

\section*{Hinweise}
\begin{itemize}
  \item Antworten Sie in ganzen Sätzen oder präzisen Stichworten.
  \item Begründen Sie Zuordnungen kurz und fachlich sauber.
  \item Skizzen und kleine Übersichten sind erlaubt.
\end{itemize}

\section*{1. Begriffe erklären}
Erklären Sie die folgenden Begriffe in \textbf{1 bis 3 eigenen Sätzen}.

\begin{longtable}{|>{\raggedright\arraybackslash}p{0.28\linewidth}|>{\raggedright\arraybackslash}p{0.64\linewidth}|}
\hline
\textbf{Begriff} & \textbf{Ihre Erklärung} \\
\hline
\term{TCP} & \\[2.1cm] \hline
\term{HTTP} & \\[2.1cm] \hline
\term{REST} & \\[2.1cm] \hline
\term{GraphQL} & \\[2.1cm] \hline
\term{gRPC} & \\[2.1cm] \hline
\term{WebSocket} & \\[2.1cm] \hline
\term{Gateway} & \\[2.1cm] \hline
\term{Proxy} & \\[2.1cm] \hline
\end{longtable}

\section*{2. Beispiele zuordnen}
Ordnen Sie jedem Beispiel den passendsten Begriff zu:

\term{TCP}, \term{HTTP}, \term{REST}, \term{GraphQL},
\term{gRPC}, \term{WebSocket}, \term{Gateway}, \term{Proxy}

\begin{enumerate}
  \item Ein Client sendet eine \texttt{GET}-Anfrage an \texttt{/api/products/5} und erhält JSON zurück.
  \item Zwei Microservices tauschen intern binär kodierte Nachrichten mit automatisch generiertem Code aus.
  \item Eine Chat-Anwendung empfängt neue Nachrichten sofort, ohne die Seite neu zu laden.
  \item Ein Client sendet eine Anfrage und fragt gezielt nur nach Name und E-Mail, nicht nach allen Feldern.
  \item Alle Anfragen laufen über einen einzigen Eintrittspunkt, der Authentifizierung und Routing übernimmt.
  \item Verlorene Datenpakete werden automatisch erneut gesendet, bis sie ankommen.
  \item Ein Webserver im Firmen-Netzwerk leitet Anfragen der Mitarbeitenden ans Internet weiter.
  \item Eine Antwort enthält den Status-Code \texttt{404 Not Found}.
\end{enumerate}

\section*{3. Protokolle und Schnittstellen vergleichen}

\begin{tabularx}{\linewidth}{|>{\raggedright\arraybackslash}p{0.30\linewidth}|>{\raggedright\arraybackslash}X|>{\raggedright\arraybackslash}X|}
\hline
\textbf{Frage} & \textbf{REST} & \textbf{GraphQL} \\
\hline
Wie viele Endpunkte hat die Schnittstelle typischerweise? & & \\[1.1cm] \hline
Wer bestimmt die Struktur der zurückgegebenen Daten? & & \\[1.1cm] \hline
In welchem Szenario ist dieser Stil besonders geeignet? & & \\[1.1cm] \hline
\end{tabularx}

\vspace{1em}

\begin{tabularx}{\linewidth}{|>{\raggedright\arraybackslash}p{0.30\linewidth}|>{\raggedright\arraybackslash}X|>{\raggedright\arraybackslash}X|}
\hline
\textbf{Frage} & \textbf{Gateway} & \textbf{Reverse Proxy} \\
\hline
Welche Hauptaufgabe hat diese Komponente? & & \\[1.1cm] \hline
Was macht sie für das System einfacher? & & \\[1.1cm] \hline
Warum sollte man beide nicht verwechseln? & & \\[1.1cm] \hline
\end{tabularx}

\section*{4. Aussagen beurteilen}
Beantworten Sie mit \textbf{richtig / falsch} und begründen Sie kurz.

\begin{enumerate}
  \item REST und HTTP sind dasselbe.
  \item GraphQL verwendet für jede Ressource einen eigenen Endpunkt.
  \item gRPC ist immer schneller als REST, unabhängig vom Anwendungsfall.
  \item Ein API-Gateway kann Authentifizierung zentral übernehmen.
  \item WebSockets eignen sich besonders gut für einfache einmalige Leseanfragen.
  \item TCP garantiert, dass Pakete in der richtigen Reihenfolge ankommen.
\end{enumerate}

\section*{5. Szenarien analysieren}
Beschreiben Sie für jedes Szenario, welche Schnittstelle oder welches Protokoll am besten passt, und begründen Sie kurz.

\begin{enumerate}
  \item Ein Entwicklungsteam baut eine öffentliche API für einen Online-Shop. Produkte sollen gelesen, erstellt und gelöscht werden können.
  \item Eine mobile App zeigt Benutzerprofil, Bestellungen und Empfehlungen an. Je nach Ansicht werden unterschiedliche Daten benötigt.
  \item Zwei interne Microservices müssen sich bei jeder Bestellung gegenseitig aufrufen. Performance und Typsicherheit sind wichtig.
  \item Eine kollaborative Textedit-Anwendung soll Änderungen sofort an alle Mitarbeitenden übermitteln.
  \item Ein Unternehmen betreibt zehn verschiedene Backend-Dienste. Alle Clients sollen nur eine einzige Adresse kennen.
\end{enumerate}

\section*{6. Schnittstellen einordnen}
Füllen Sie die Tabelle aus.

\begin{tabularx}{\linewidth}{|>{\raggedright\arraybackslash}p{0.18\linewidth}|>{\raggedright\arraybackslash}X|>{\raggedright\arraybackslash}X|>{\raggedright\arraybackslash}X|}
\hline
\textbf{Stil} & \textbf{Datenformat} & \textbf{Kommunikations\-muster} & \textbf{Typischer Einsatz} \\
\hline
REST & & & \\[0.9cm] \hline
GraphQL & & & \\[0.9cm] \hline
gRPC & & & \\[0.9cm] \hline
WebSocket & & & \\[0.9cm] \hline
\end{tabularx}

\section*{7. Mini-Entwurf}
Entwerfen Sie für eine einfache Lieferdienst-Anwendung eine grobe Schnittstellenübersicht.

Funktionen:
\begin{itemize}
  \item Benutzer melden sich an und verwalten ihre Bestellungen.
  \item Eine mobile App soll genau die Daten laden, die sie gerade braucht.
  \item Interne Dienste (Lager, Bezahlung, Versand) müssen schnell miteinander kommunizieren.
  \item Fahrer erhalten Live-Updates zu neuen Aufträgen.
  \item Alle externen Anfragen sollen über einen einzigen Zugangspunkt laufen.
\end{itemize}

\begin{enumerate}
  \item Welche Schnittstellenstile oder Protokolle würden Sie einsetzen?
  \item Begründen Sie Ihre Wahl für jeden Bereich kurz.
  \item Welche Komponente steuert den Zugang von aussen?
\end{enumerate}

\section*{8. Reflexion}
Beantworten Sie die Fragen kurz.

\begin{enumerate}
  \item Warum ist es wichtig, dass Schnittstellen klar definiert und dokumentiert sind?
  \item Warum sollte man nicht immer die neueste oder komplexeste Schnittstellentechnologie wählen?
  \item Was ist der Unterschied zwischen einem Protokoll und einem Architekturstil wie REST?
\end{enumerate}

\end{document}
